\documentclass[10pt,a4paper]{amsart}
\usepackage{amsmath,amssymb,amsthm,graphicx}
\usepackage{verbatim}
\newenvironment{reference}{\comment}{\endcomment}
%\newenvironment{reference}{}{}
\newenvironment{slogan}{\comment}{\endcomment}
\newenvironment{history}{\comment}{\endcomment}

% For commutative diagrams we use Xy-pic
\usepackage[all]{xy}

% We use 2cell for 2-commutative diagrams.
\xyoption{2cell}
% Theorem environments.
%
\theoremstyle{plain}
\newtheorem{theorem}[subsection]{Theorem}
\newtheorem{proposition}[subsection]{Proposition}
\newtheorem{lemma}[subsection]{Lemma}

\theoremstyle{definition}
\newtheorem{definition}[subsection]{Definition}
\newtheorem{example}[subsection]{Example}
\newtheorem{exercise}[subsection]{Exercise}
\newtheorem{situation}[subsection]{Situation}

\theoremstyle{remark}
\newtheorem{remark}[subsection]{Remark}
\newtheorem{remarks}[subsection]{Remarks}

\numberwithin{equation}{subsection}

% Macros
%
\def\lim{\mathop{\mathrm{lim}}\nolimits}
\def\colim{\mathop{\mathrm{colim}}\nolimits}
\def\Spec{\mathop{\mathrm{Spec}}}
\def\Hom{\mathop{\mathrm{Hom}}\nolimits}
\def\Ext{\mathop{\mathrm{Ext}}\nolimits}
\def\SheafHom{\mathop{\mathcal{H}\!\mathit{om}}\nolimits}
\def\SheafExt{\mathop{\mathcal{E}\!\mathit{xt}}\nolimits}
\def\Sch{\mathit{Sch}}
\def\Mor{\mathop{\mathrm{Mor}}\nolimits}
\def\Ob{\mathop{\mathrm{Ob}}\nolimits}
\def\Sh{\mathop{\mathit{Sh}}\nolimits}
\def\NL{\mathop{N\!L}\nolimits}
\def\CH{\mathop{\mathrm{CH}}\nolimits}
\def\proetale{{pro\text{-}\acute{e}tale}}
\def\etale{{\acute{e}tale}}
\def\QCoh{\mathit{QCoh}}
\def\Ker{\mathop{\mathrm{Ker}}}
\def\Im{\mathop{\mathrm{Im}}}
\def\Coker{\mathop{\mathrm{Coker}}}
\def\Coim{\mathop{\mathrm{Coim}}}

% Boxtimes
%
\DeclareMathSymbol{\boxtimes}{\mathbin}{AMSa}{"02}

%
% Macros for moduli stacks/spaces
%
\def\QCohstack{\mathcal{QC}\!\mathit{oh}}
\def\Cohstack{\mathcal{C}\!\mathit{oh}}
\def\Spacesstack{\mathcal{S}\!\mathit{paces}}
\def\Quotfunctor{\mathrm{Quot}}
\def\Hilbfunctor{\mathrm{Hilb}}
\def\Curvesstack{\mathcal{C}\!\mathit{urves}}
\def\Polarizedstack{\mathcal{P}\!\mathit{olarized}}
\def\Complexesstack{\mathcal{C}\!\mathit{omplexes}}
% \Pic is the operator that assigns to X its picard group, usage \Pic(X)
% \Picardstack_{X/B} denotes the Picard stack of X over B
% \Picardfunctor_{X/B} denotes the Picard functor of X over B
\def\Pic{\mathop{\mathrm{Pic}}\nolimits}
\def\Picardstack{\mathcal{P}\!\mathit{ic}}
\def\Picardfunctor{\mathrm{Pic}}
\def\Deformationcategory{\mathcal{D}\!\mathit{ef}}
\usepackage{cleveref}
\crefname{originalproposition}{original proposition}{original propositions}
\Crefname{originalproposition}{Original proposition}{Original propositions}
\makeatletter
\newlabel{algebra-proposition-ubiquity-nagata}{{0335}{source}{}{sourcepointer-1}{}}
\newlabel{algebra-proposition-ubiquity-nagata@cref}{{[originalproposition][0][]0335}{[1][0][]source}{}{}{}}
\makeatother
\crefname{originallemma}{original lemma}{original lemmas}
\Crefname{originallemma}{Original lemma}{Original lemmas}
\makeatletter
\newlabel{dualizing-lemma-ubiquity-dualizing}{{0BFR}{source}{}{sourcepointer-2}{}}
\newlabel{dualizing-lemma-ubiquity-dualizing@cref}{{[originallemma][0][]0BFR}{[1][0][]source}{}{}{}}
\makeatother
\crefname{originalreference}{original reference}{original references}
\Crefname{originalreference}{Original reference}{Original references}
\makeatletter
\newlabel{lemma-dominate-by-blowing-up-in-points}{{0AHI}{source}{}{sourcepointer-3}{}}
\newlabel{lemma-dominate-by-blowing-up-in-points@cref}{{[originalreference][0][]0AHI}{[1][0][]source}{}{}{}}
\makeatother
\crefname{originalreference}{original reference}{original references}
\Crefname{originalreference}{Original reference}{Original references}
\makeatletter
\newlabel{lemma-dominate-by-normalized-blowing-up}{{0BBT}{source}{}{sourcepointer-4}{}}
\newlabel{lemma-dominate-by-normalized-blowing-up@cref}{{[originalreference][0][]0BBT}{[1][0][]source}{}{}{}}
\makeatother
\crefname{originalreference}{original reference}{original references}
\Crefname{originalreference}{Original reference}{Original references}
\makeatletter
\newlabel{lemma-equivalence-sequence-blowups}{{0BFY}{source}{}{sourcepointer-5}{}}
\newlabel{lemma-equivalence-sequence-blowups@cref}{{[originalreference][0][]0BFY}{[1][0][]source}{}{}{}}
\makeatother
\crefname{originalreference}{original reference}{original references}
\Crefname{originalreference}{Original reference}{Original references}
\makeatletter
\newlabel{lemma-normalized-blowup-completion}{{0BGA}{source}{}{sourcepointer-6}{}}
\newlabel{lemma-normalized-blowup-completion@cref}{{[originalreference][0][]0BGA}{[1][0][]source}{}{}{}}
\makeatother
\crefname{originalreference}{original reference}{original references}
\Crefname{originalreference}{Original reference}{Original references}
\makeatletter
\newlabel{lemma-port-regularity-to-completion}{{0BG6}{source}{}{sourcepointer-7}{}}
\newlabel{lemma-port-regularity-to-completion@cref}{{[originalreference][0][]0BG6}{[1][0][]source}{}{}{}}
\makeatother
\crefname{originalreference}{original reference}{original references}
\Crefname{originalreference}{Original reference}{Original references}
\makeatletter
\newlabel{lemma-rational-propagates}{{0BG0}{source}{}{sourcepointer-8}{}}
\newlabel{lemma-rational-propagates@cref}{{[originalreference][0][]0BG0}{[1][0][]source}{}{}{}}
\makeatother
\crefname{originalreference}{original reference}{original references}
\Crefname{originalreference}{Original reference}{Original references}
\makeatletter
\newlabel{lemma-rational-to-gorenstein}{{0BBV}{source}{}{sourcepointer-9}{}}
\newlabel{lemma-rational-to-gorenstein@cref}{{[originalreference][0][]0BBV}{[1][0][]source}{}{}{}}
\makeatother
\crefname{originalreference}{original reference}{original references}
\Crefname{originalreference}{Original reference}{Original references}
\makeatletter
\newlabel{lemma-regular-alteration-implies-local}{{0BGI}{source}{}{sourcepointer-10}{}}
\newlabel{lemma-regular-alteration-implies-local@cref}{{[originalreference][0][]0BGI}{[1][0][]source}{}{}{}}
\makeatother
\crefname{originalreference}{original reference}{original references}
\Crefname{originalreference}{Original reference}{Original references}
\makeatletter
\newlabel{lemma-regular-rational}{{0B4Q}{source}{}{sourcepointer-11}{}}
\newlabel{lemma-regular-rational@cref}{{[originalreference][0][]0B4Q}{[1][0][]source}{}{}{}}
\makeatother
\crefname{originalreference}{original reference}{original references}
\Crefname{originalreference}{Original reference}{Original references}
\makeatletter
\newlabel{lemma-resolve-rational-double-points}{{0BGE}{source}{}{sourcepointer-12}{}}
\newlabel{lemma-resolve-rational-double-points@cref}{{[originalreference][0][]0BGE}{[1][0][]source}{}{}{}}
\makeatother
\crefname{originallemma}{original lemma}{original lemmas}
\Crefname{originallemma}{Original lemma}{Original lemmas}
\makeatletter
\newlabel{more-algebra-lemma-normal-goes-up}{{0BFK}{source}{}{sourcepointer-13}{}}
\newlabel{more-algebra-lemma-normal-goes-up@cref}{{[originallemma][0][]0BFK}{[1][0][]source}{}{}{}}
\makeatother
\crefname{originalproposition}{original proposition}{original propositions}
\Crefname{originalproposition}{Original proposition}{Original propositions}
\makeatletter
\newlabel{more-algebra-proposition-ubiquity-excellent}{{07QW}{source}{}{sourcepointer-14}{}}
\newlabel{more-algebra-proposition-ubiquity-excellent@cref}{{[originalproposition][0][]07QW}{[1][0][]source}{}{}{}}
\makeatother
\crefname{originallemma}{original lemma}{original lemmas}
\Crefname{originallemma}{Original lemma}{Original lemmas}
\makeatletter
\newlabel{varieties-lemma-modification-normal-iso-over-codimension-1}{{0BFP}{source}{}{sourcepointer-15}{}}
\newlabel{varieties-lemma-modification-normal-iso-over-codimension-1@cref}{{[originallemma][0][]0BFP}{[1][0][]source}{}{}{}}
\makeatother
\title{A resolution of a normal two dimensional local ring yields a resolution by normalized blowups}
\author{The Stacks Project Authors}
\date{stacks-project:resolve.tex:local-snapshot-20261012; GNU FDL 1.2 or later}
\begin{document}
\maketitle
\noindent\textit{Editorial note.} All delivered mathematical text is the author's own, copied from the pinned Stacks Project chapter file resolve.tex; no statement or proof step was rewritten, reordered or removed.\par\medskip
\noindent\textit{Editorial note.} The counted proof is the author's complete argument: the reduction to the complete case through the transfer of normality and regularity along completion and the descent of sequences of normalized blowups, the domination of the given resolution first by normalized blowups and then by blowups in closed points yielding finitely many singular points, the propagation of rationality followed by the passage to Gorenstein local rings, and the final resolution of rational double points.\par\medskip
\noindent\textit{Editorial note.} The label lemma-existence-implies-existence-by-normalized-blowing-ups names the counted statement, and the proof environment immediately following it in the source is delivered verbatim as its proof.\par\medskip
\noindent\textit{Editorial note.} The counted statement is carried by the author's own lemma environment in the chapter source.\par\medskip
\noindent\textit{Editorial note.} Pre-open gate: python operations/agent-authoring/tools/precard\_gate.py --source resolve.tex --label lemma-existence-implies-existence-by-normalized-blowing-ups --start 3949 --end 4048 --local-id 28073 returned RESULT=OK (L28073 free, occupancy FREE, public interval clear). Pre-assembly tag self-check: canonical tag 0BGM, not present in reference/public-items.jsonl.\par\medskip
\noindent\textit{Editorial note.} The statement and proof are copied from the pinned chapter file; the cited inputs are kept as external pointers with their canonical Stacks tags.\par\medskip
\noindent\textit{Allowed standard prerequisites (editorial).}\begin{itemize}
\item Normality and regularity pass to and from the completion. Same chapter, the lemma that the completion of a normal local ring of a regular alteration is normal, and the lemma transferring regularity of a resolution along the completion. Supplies the reduction to the complete case and the transfer of regularity between a sequence and its completion.
\item Descending a sequence of normalized blowups from the completion. Same chapter, the lemma that a sequence of normalized blowups of the completion comes from a sequence over the original ring. Supplies the passage from sequences over the completion back to sequences over the ring.
\item Domination of a resolution by normalized blowups at closed points. Same chapter, the lemma dominating a proper birational morphism of a two dimensional Noetherian Nagata scheme by normalized blowups at closed points, and the lemma dominating a modification by blowups in points. Supply the sequence of normalized blowups dominating the given resolution and the further sequence of point blowups dominating it.
\item Propagation of rational singularities and the resolution of rational double points. Same chapter, the lemma that regularity defines a rational singularity, the lemma propagating the rational singularity property, the lemma producing finitely many blowups making the local rings Gorenstein, and the lemma resolving rational double points. Supply the intermediate surface with Gorenstein local rings and rational double point singularities, and its resolution.
\end{itemize}\medskip
\ifdefined\hypertarget\hypertarget{sourcepointer-1}{}\fi
\noindent\textit{Source reference.} Original proposition 0335: sources/stacks-chow/tags.txt, source lines 4044--4044. Prerequisite pointer invoked at proof line(s) 3972; it names a standard fact of a sibling chapter that the counted argument uses, and every step of the counted proof built on it is retained verbatim in the delivered body.\par\medskip
\ifdefined\hypertarget\hypertarget{sourcepointer-2}{}\fi
\noindent\textit{Source reference.} Original lemma 0BFR: sources/stacks-chow/tags.txt, source lines 14394--14394. Prerequisite pointer invoked at proof line(s) 3976; it names a standard fact of a sibling chapter that the counted argument uses, and every step of the counted proof built on it is retained verbatim in the delivered body.\par\medskip
\ifdefined\hypertarget\hypertarget{sourcepointer-3}{}\fi
\noindent\textit{Source reference.} Original reference 0AHI: sources/stacks-resolve/pinned-original-source.tex, source lines 1146--1146. Prerequisite pointer invoked at proof line(s) 3992; it names a standard fact of the same chapter that the counted argument uses, and every step of the counted proof built on it is retained verbatim in the delivered body.\par\medskip
\ifdefined\hypertarget\hypertarget{sourcepointer-4}{}\fi
\noindent\textit{Source reference.} Original reference 0BBT: sources/stacks-resolve/pinned-original-source.tex, source lines 1277--1277. Prerequisite pointer invoked at proof line(s) 3989; it names a standard fact of the same chapter that the counted argument uses, and every step of the counted proof built on it is retained verbatim in the delivered body.\par\medskip
\ifdefined\hypertarget\hypertarget{sourcepointer-5}{}\fi
\noindent\textit{Source reference.} Original reference 0BFY: sources/stacks-resolve/pinned-original-source.tex, source lines 1495--1495. Prerequisite pointer invoked at proof line(s) 4026; it names a standard fact of the same chapter that the counted argument uses, and every step of the counted proof built on it is retained verbatim in the delivered body.\par\medskip
\ifdefined\hypertarget\hypertarget{sourcepointer-6}{}\fi
\noindent\textit{Source reference.} Original reference 0BGA: sources/stacks-resolve/pinned-original-source.tex, source lines 3236--3236. Prerequisite pointer invoked at proof line(s) 3962; it names a standard fact of the same chapter that the counted argument uses, and every step of the counted proof built on it is retained verbatim in the delivered body.\par\medskip
\ifdefined\hypertarget\hypertarget{sourcepointer-7}{}\fi
\noindent\textit{Source reference.} Original reference 0BG6: sources/stacks-resolve/pinned-original-source.tex, source lines 2997--2997. Prerequisite pointer invoked at proof line(s) 3960; it names a standard fact of the same chapter that the counted argument uses, and every step of the counted proof built on it is retained verbatim in the delivered body.\par\medskip
\ifdefined\hypertarget\hypertarget{sourcepointer-8}{}\fi
\noindent\textit{Source reference.} Original reference 0BG0: sources/stacks-resolve/pinned-original-source.tex, source lines 2035--2035. Prerequisite pointer invoked at proof line(s) 4016; it names a standard fact of the same chapter that the counted argument uses, and every step of the counted proof built on it is retained verbatim in the delivered body.\par\medskip
\ifdefined\hypertarget\hypertarget{sourcepointer-9}{}\fi
\noindent\textit{Source reference.} Original reference 0BBV: sources/stacks-resolve/pinned-original-source.tex, source lines 2688--2688. Prerequisite pointer invoked at proof line(s) 4019; it names a standard fact of the same chapter that the counted argument uses, and every step of the counted proof built on it is retained verbatim in the delivered body.\par\medskip
\ifdefined\hypertarget\hypertarget{sourcepointer-10}{}\fi
\noindent\textit{Source reference.} Original reference 0BGI: sources/stacks-resolve/pinned-original-source.tex, source lines 3827--3827. Prerequisite pointer invoked at proof line(s) 3958; it names a standard fact of the same chapter that the counted argument uses, and every step of the counted proof built on it is retained verbatim in the delivered body.\par\medskip
\ifdefined\hypertarget\hypertarget{sourcepointer-11}{}\fi
\noindent\textit{Source reference.} Original reference 0B4Q: sources/stacks-resolve/pinned-original-source.tex, source lines 2205--2205. Prerequisite pointer invoked at proof line(s) 4015; it names a standard fact of the same chapter that the counted argument uses, and every step of the counted proof built on it is retained verbatim in the delivered body.\par\medskip
\ifdefined\hypertarget\hypertarget{sourcepointer-12}{}\fi
\noindent\textit{Source reference.} Original reference 0BGE: sources/stacks-resolve/pinned-original-source.tex, source lines 3716--3716. Prerequisite pointer invoked at proof line(s) 4046; it names a standard fact of the same chapter that the counted argument uses, and every step of the counted proof built on it is retained verbatim in the delivered body.\par\medskip
\ifdefined\hypertarget\hypertarget{sourcepointer-13}{}\fi
\noindent\textit{Source reference.} Original lemma 0BFK: sources/stacks-chow/tags.txt, source lines 14388--14388. Prerequisite pointer invoked at proof line(s) 3980; it names a standard fact of a sibling chapter that the counted argument uses, and every step of the counted proof built on it is retained verbatim in the delivered body.\par\medskip
\ifdefined\hypertarget\hypertarget{sourcepointer-14}{}\fi
\noindent\textit{Source reference.} Original proposition 07QW: sources/stacks-chow/tags.txt, source lines 9846--9846. Prerequisite pointer invoked at proof line(s) 3974; it names a standard fact of a sibling chapter that the counted argument uses, and every step of the counted proof built on it is retained verbatim in the delivered body.\par\medskip
\ifdefined\hypertarget\hypertarget{sourcepointer-15}{}\fi
\noindent\textit{Source reference.} Original lemma 0BFP: sources/stacks-chow/tags.txt, source lines 14392--14392. Prerequisite pointer invoked at proof line(s) 4010; it names a standard fact of a sibling chapter that the counted argument uses, and every step of the counted proof built on it is retained verbatim in the delivered body.\par\medskip
\begin{lemma}
\label{lemma-existence-implies-existence-by-normalized-blowing-ups}
Let $(A, \mathfrak m, \kappa)$ be a Noetherian local ring.
Assume $A$ is normal and has dimension $2$.
If $\Spec(A)$ has a resolution of singularities,
then $\Spec(A)$ has a resolution by normalized blowups.
\end{lemma}

\begin{proof}
By Lemma \ref{lemma-regular-alteration-implies-local}
the completion $A^\wedge$ of $A$ is normal.
By Lemma \ref{lemma-port-regularity-to-completion} we see
that $\Spec(A^\wedge)$ has a resolution.
By Lemma \ref{lemma-normalized-blowup-completion}
any sequence $Y_n \to Y_{n - 1} \to \ldots \to \Spec(A^\wedge)$
of normalized blowups of comes from a sequence of normalized
blowups $X_n \to \ldots \to \Spec(A)$. Moreover if $Y_n$ is
regular, then $X_n$ is regular by
Lemma \ref{lemma-port-regularity-to-completion}.
Thus it suffices to prove the lemma in case $A$ is complete.

\medskip\noindent
Assume in addition $A$ is a complete. We will use that $A$ is Nagata
(Algebra, Proposition \ref{algebra-proposition-ubiquity-nagata}),
excellent (More on Algebra, Proposition
\ref{more-algebra-proposition-ubiquity-excellent}),
and has a dualizing complex
(Dualizing Complexes, Lemma \ref{dualizing-lemma-ubiquity-dualizing}).
Moreover, the same is true for any ring essentially of finite type over $A$.
If $B$ is a excellent local normal domain, then the completion
$B^\wedge$ is normal (as $B \to B^\wedge$ is regular and
More on Algebra, Lemma \ref{more-algebra-lemma-normal-goes-up} applies).
We will use this without further mention in the rest of the proof.

\medskip\noindent
Let $X \to \Spec(A)$ be a resolution of singularities.
Choose a sequence of normalized blowing ups
$$
Y_n \to Y_{n - 1} \to \ldots \to Y_1 \to \Spec(A)
$$
dominating $X$ (Lemma \ref{lemma-dominate-by-normalized-blowing-up}).
The morphism $Y_n \to X$ is an isomorphism away from
finitely many points of $X$.
Hence we can apply Lemma \ref{lemma-dominate-by-blowing-up-in-points}
to find a sequence of blowing ups
$$
X_m \to X_{m - 1} \to \ldots \to X
$$
in closed points such that $X_m$ dominates $Y_n$. Diagram
$$
\xymatrix{
& Y_n \ar[rd] \ar[rr] & & \Spec(A) \\
X_m \ar[rr] \ar[ru] & & X \ar[ru]
}
$$
To prove the lemma it suffices to show that a finite number of normalized
blowups of $Y_n$ produce a regular scheme. By our diagram above we see that
$Y_n$ has a resolution (namely $X_m$). As $Y_n$ is a normal surface
this implies that $Y_n$ has at most finitely many singularities
$y_1, \ldots, y_t$ (because $X_m \to Y_n$ is an isomorphism away from
the fibres of dimension $1$, see Varieties, Lemma
\ref{varieties-lemma-modification-normal-iso-over-codimension-1}).

\medskip\noindent
Let $x_a \in X$ be the image of $y_a$. Then $\mathcal{O}_{X, x_a}$
is regular and hence defines a rational singularity
(Lemma \ref{lemma-regular-rational}).
Apply Lemma \ref{lemma-rational-propagates} to
$\mathcal{O}_{X, x_a} \to \mathcal{O}_{Y_n, y_a}$
to see that $\mathcal{O}_{Y_n, y_a}$ defines a
rational singularity. By Lemma \ref{lemma-rational-to-gorenstein}
there exists a finite sequence of blowups in singular closed points
$$
Y_{a, n_a} \to Y_{a, n_a - 1} \to \ldots \to \Spec(\mathcal{O}_{Y_n, y_a})
$$
such that $Y_{a, n_a}$ is Gorenstein, i.e., has an
invertible dualizing module. By (the essentially trivial)
Lemma \ref{lemma-equivalence-sequence-blowups}
with $n' = \sum n_a$ these sequences correspond to a sequence of
blowups
$$
Y_{n + n'} \to Y_{n + n' - 1} \to \ldots \to Y_n
$$
such that $Y_{n + n'}$ is normal and
the local rings of $Y_{n + n'}$ are Gorenstein. Using the references
given above we can dominate $Y_{n + n'}$ by a sequence of blowups
$X_{m + m'} \to \ldots \to X_m$ dominating $Y_{n + n'}$ as in the following
$$
\xymatrix{
 & Y_{n + n'} \ar[rr] & & Y_n \ar[rd] \ar[rr] & & \Spec(A) \\
X_{m + m'} \ar[ru] \ar[rr] & & X_m \ar[rr] \ar[ru] & & X \ar[ru]
}
$$
Thus again $Y_{n + n'}$ has a finite number of singular points
$y'_1, \ldots, y'_s$, but this time the singularities are
rational double points, more precisely, the local rings
$\mathcal{O}_{Y_{n + n'}, y'_b}$ are as in
Lemma \ref{lemma-resolve-rational-double-points}.
Arguing exactly as above we conclude that the lemma is true.
\end{proof}

\end{document}
